\section{Introduction}
\label{sec:jesper-introduction}
\contributionauthor{Jesper Eggers}
Bomberman makes a useful test of reinforcement learning because an action
can create both an opportunity and a delayed threat to its own actor.
Bombs clear terrain and eliminate opponents, but each placement creates
a delayed hazard that the agent must anticipate. Navigation also changes
as crates disappear, while opponents compete for coins and alter the
available routes. A high shaped reward or successful play without
opponents therefore need not predict competitive strength.

The six discrete actions suit value-based learning. Tabular Q-learning
provides a conceptual starting point, but enumerating configurations of
terrain, bombs, and opponents quickly becomes impractical. Neural action
values offer generalization across configurations; direct policy optimization
would be another possible approach. We specialize the value-based approach
through engineered observations, dueling networks, reward shaping, and masks.

Our project implements two agents based on Double DQN and dueling
networks. RUEHL uses a compact learned representation and comparatively
simple action constraints. The final agent combines a larger representation
with explicit temporal survival searches; its name does not denote the
submitted player, which was RUEHL. The central experimental question
is how learned action preferences and safety guidance contribute to their
deployed behavior. Alongside the existing development results, we test the
masks against engine-derived continuations and evaluate frozen final-agent
weights in a controlled action-selection experiment. This separates
constraint quality, learned ranking, and competitive score without adding
another architecture or training campaign.

We distinguish engine-conditional escape feasibility, learned choices
within the constraints, and competitive outcomes. Temporal guidance and
learned ranking both add score in the tested deployment configurations,
while learned selection increases self-destruction relative to masked
random play. A replayed opponent collision explains why conditional
reachability remains weaker than competitive safety.
